\begin{frame}{Árvore de segmentos - soma}
    \metroset{block=fill}

    Seja $[l, r]$ o intervalo que queremos calcular. Supondo que a árvore de segmentos já está montada, o algoritmo para calcular a soma de $l$ até $r$ é:
    \begin{itemize}
        \item Começamos na raiz da árvore, que contém a soma do intervalo $[1, n]$. O algoritmo será recursivo em função do nó.
        \item Considere que estamos no nó correspondente ao intervalo $[lx, rx]$. Se $lx > rx$ então retornamos $0$. Caso contrário, consideramos dois casos:
        \item Caso 1: $[lx, rx] \not \subset [l, r]$. Isso significa que existem elementos no intervalo $[lx, rx]$ que não queremos na soma. Chamamos recursivamente para os intervalos $[lx, mid]$ e $[mid+1, rx]$, onde $mid = (lx+rx)/2$, e retornamos a soma dos retornos dessas chamadas.
        \item Caso 2: $[lx, rx] \subset [l, r]$. Neste caso, retornamos a soma do nó atual.
    \end{itemize}

    % \begin{block}{Implementação da busca binária}
    %     \footnotesize
    %     \inputminted{cpp}{codigos/busca.binaria.cpp}    
    % \end{block}

\end{frame}
